\textbf{Phần 1. Trắc nghiệm lựa chọn}
\setcounterpageref{ex}{0}

\begin{ex} %[2D2V13]
Một bến xe có $17$ chiếc ô tô. Buổi sáng có $8$ chiếc rời bến, buổi chiều có thêm $5$ chiếc vào bến. Hỏi lúc này bến xe có bao nhiêu chiếc ô tô?
\choice
{$13$ chiếc ô tô}
{$15$ chiếc ô tô}
{\True $14$ chiếc ô tô}
{$12$ chiếc ô tô}
\loigiai{
Số ô tô lúc này: $17 - 8 + 5 = 9 + 5 = 14$ (chiếc ô tô).
\textbf{Đáp án của câu hỏi là C}
}
\end{ex}

\begin{ex} %[2D2V13]
Mẹ mua về một túi mận. Sau khi cả nhà ăn hết $9$ quả thì trong túi còn lại $8$ quả. Hỏi ban đầu túi mận có bao nhiêu quả?
\choice
{\True $17$ quả}
{$16$ quả}
{$18$ quả}
{$15$ quả}
\loigiai{
Ban đầu túi mận có: $9 + 8 = 17$ (quả).
\textbf{Đáp án của câu hỏi là A}
}
\end{ex}

\begin{ex} %[2D2V13]
Một đàn vịt có $14$ con. Người ta thả thêm vào đàn một số con vịt thì đàn vịt có tất cả $20$ con. Hỏi người ta đã thả thêm bao nhiêu con vịt?
\choice
{$5$ con vịt}
{$7$ con vịt}
{$8$ con vịt}
{\True $6$ con vịt}
\loigiai{
Số vịt thả thêm là: $20 - 14 = 6$ (con vịt).
\textbf{Đáp án của câu hỏi là D}
}
\end{ex}

\begin{ex} %[2D2V13]
Bình có $16$ viên bi. Bình cho Nam $5$ viên bi rồi lại được bố cho thêm $4$ viên bi. Hỏi lúc này Bình có bao nhiêu viên bi?
\choice
{$14$ viên bi}
{\True $15$ viên bi}
{$16$ viên bi}
{$13$ viên bi}
\loigiai{
Số bi của Bình lúc này là: $16 - 5 + 4 = 11 + 4 = 15$ (viên bi).
\textbf{Đáp án của câu hỏi là B}
}
\end{ex}

\begin{ex} %[2D2V13]
Trong một cuộc thi chạy có $15$ vận động viên. Có $6$ vận động viên đã về đích. Hỏi còn bao nhiêu vận động viên đang trên đường chạy?
\choice
{$8$ người}
{$10$ người}
{\True $9$ người}
{$7$ người}
\loigiai{
Số vận động viên còn lại: $15 - 6 = 9$ (người).
\textbf{Đáp án của câu hỏi là C}
}
\end{ex}

\begin{ex} %[2D2V13]
Một cuộn dây dài $18\text{ m}$. Bác thợ cắt đi $7\text{ m}$ để buộc hàng, sau đó cắt tiếp $4\text{ m}$ để làm dây phơi. Hỏi cuộn dây còn lại bao nhiêu mét?
\choice
{\True $7\text{ m}$}
{$6\text{ m}$}
{$8\text{ m}$}
{$9\text{ m}$}
\loigiai{
Tổng số mét dây đã cắt: $7 + 4 = 11\text{ (m)}$.\nCuộn dây còn lại: $18 - 11 = 7\text{ (m)}$.
\textbf{Đáp án của câu hỏi là A}
}
\end{ex}

\begin{ex} %[2D2V13]
Lớp 2B có $18$ bạn tham gia văn nghệ. Nếu có thêm $2$ bạn nam nữa thì số bạn tham gia văn nghệ là
\choice
{$19$ bạn}
{$16$ bạn}
{$21$ bạn}
{\True $20$ bạn}
\loigiai{
$18 + 2 = 20$ (bạn).
\textbf{Đáp án của câu hỏi là D}
}
\end{ex}

\begin{ex} %[2D2V13]
Một rổ táo có $13$ quả. Sau khi bỏ thêm vào rổ $6$ quả nữa rồi lấy ra $8$ quả thì trong rổ còn lại bao nhiêu quả?
\choice
{$12$ quả}
{\True $11$ quả}
{$10$ quả}
{$13$ quả}
\loigiai{
$13 + 6 - 8 = 19 - 8 = 11$ (quả).
\textbf{Đáp án của câu hỏi là B}
}
\end{ex}

\begin{ex} %[2D2V13]
Tìm một số biết rằng lấy số đó bớt đi $7$ rồi thêm vào $9$ thì được $18$.
\choice
{$15$}
{$17$}
{\True $16$}
{$14$}
\loigiai{
Tính ngược từ cuối: $18 - 9 + 7 = 9 + 7 = 16$.
\textbf{Đáp án của câu hỏi là C}
}
\end{ex}

\begin{ex} %[2D2V13]
Một bể cá có $15$ con cá vàng. Người ta vớt ra một số con cá đem tặng thì trong bể còn lại $8$ con. Hỏi người ta đã vớt ra bao nhiêu con cá?
\choice
{\True $7$ con cá}
{$6$ con cá}
{$8$ con cá}
{$9$ con cá}
\loigiai{
Số cá vớt ra là: $15 - 8 = 7$ (con cá).
\textbf{Đáp án của câu hỏi là A}
}
\end{ex}

\begin{ex} %[2D2V13]
Mai hái được $8$ quả dâu tây, Lan hái được $9$ quả dâu tây. Hai bạn cùng ăn hết $5$ quả. Hỏi hai bạn còn lại bao nhiêu quả dâu tây?
\choice
{$11$ quả}
{$13$ quả}
{$14$ quả}
{\True $12$ quả}
\loigiai{
Tổng số dâu tây: $8 + 9 = 17$ (quả).\nCòn lại: $17 - 5 = 12$ (quả).
\textbf{Đáp án của câu hỏi là D}
}
\end{ex}

\begin{ex} %[2D2V13]
Một cửa hàng buổi sáng có $19$ chiếc mũ bảo hiểm, buổi chiều bán được $7$ chiếc. Hỏi cửa hàng còn lại bao nhiêu chiếc mũ?
\choice
{$11$ chiếc}
{\True $12$ chiếc}
{$13$ chiếc}
{$10$ chiếc}
\loigiai{
$19 - 7 = 12$ (chiếc).
\textbf{Đáp án của câu hỏi là B}
}
\end{ex}

\textbf{Phần 2. Trắc nghiệm đúng sai}
\setcounterpageref{ex}{0}

\begin{ex} %[2D2V13]
Xét tính đúng sai của các bài toán thêm, bớt:
\choiceTF
{\True Bớt đi ở số ban đầu một số lượng thì số mới sẽ nhỏ hơn số ban đầu.}
{Thêm vào số ban đầu 5 đơn vị rồi bớt đi 5 đơn vị thì số mới lớn hơn số ban đầu.}
{\True Một xe có 14 chỗ ngồi, đã có 9 khách ngồi thì còn lại 5 chỗ trống.}
{\True Lấy một số cộng với 0 thì vẫn bằng chính số đó.}
\loigiai{
\begin{itemize}
    \item Ý a) Đúng: Phép trừ luôn làm giảm giá trị của số tự nhiên khác 0.
    \item Ý b) Sai: Số mới bằng đúng số ban đầu: $+ 5 - 5 = 0$.
    \item Ý c) Đúng: $14 - 9 = 5$ chỗ trống. Đúng.
    \item Ý d) Đúng: Tính chất cộng với số 0.
\end{itemize}
\begin{itemize}
    \item a) \textbf{Đúng}.
    \item b) \textbf{Sai}.
    \item c) \textbf{Đúng}.
    \item d) \textbf{Đúng}.
\end{itemize}
}
\end{ex}

\begin{ex} %[2D2V13]
Trên đĩa có 16 cái bánh. Nam ăn 3 cái, Mai ăn 4 cái:
\choiceTF
{\True Cả hai bạn ăn tất cả 7 cái bánh.}
{\True Trên đĩa còn lại 9 cái bánh.}
{Số bánh còn lại ít hơn số bánh đã ăn.}
{Nếu mẹ cho thêm vào đĩa 2 cái bánh thì trên đĩa có tất cả 10 cái bánh.}
\loigiai{
\begin{itemize}
    \item Ý a) Đúng: $3 + 4 = 7$ cái bánh.
    \item Ý b) Đúng: $16 - 7 = 9$ cái bánh.
    \item Ý c) Sai: Còn 9 cái bánh, đã ăn 7 cái, vậy số còn lại nhiều hơn số đã ăn.
    \item Ý d) Sai: Số bánh lúc đó là: $9 + 2 = 11$ cái (không phải 10 cái).
\end{itemize}
\begin{itemize}
    \item a) \textbf{Đúng}.
    \item b) \textbf{Đúng}.
    \item c) \textbf{Sai}.
    \item d) \textbf{Sai}.
\end{itemize}
}
\end{ex}

\textbf{Phần 3. Trả lời ngắn}
\setcounterpageref{ex}{0}

\begin{ex} %[2D2V13]
Một cây táo có 17 quả. Sau một đêm gió bão có 8 quả táo bị rụng. Hỏi trên cây còn lại bao nhiêu quả táo?
\shortans{9}
\loigiai{
\textbf{Đáp số:} $9$\\
$17 - 8 = 9$ quả.
}
\end{ex}

\begin{ex} %[2D2V13]
Một bác nông dân gặt lúa, buổi sáng gặt được 8 bó lúa, buổi chiều gặt thêm 9 bó lúa nữa. Bác chở về nhà 11 bó. Hỏi còn lại bao nhiêu bó lúa ở ngoài đồng?
\shortans{6}
\loigiai{
\textbf{Đáp số:} $6$\\
Tổng số bó lúa: $8 + 9 = 17$. Còn lại: $17 - 11 = 6$ bó.
}
\end{ex}

\begin{ex} %[2D2V13]
Tìm một số biết rằng lấy số đó thêm 6 đơn vị rồi bớt đi 9 đơn vị thì được 10.
\shortans{13}
\loigiai{
\textbf{Đáp số:} $13$\\
Số đó là: $10 + 9 - 6 = 13$.
}
\end{ex}

\begin{ex} %[2D2V13]
Huy có 15 viên bi, Huy cho bạn 7 viên bi rồi lại được tặng thêm 5 viên bi. Hỏi Huy có tất cả bao nhiêu viên bi?
\shortans{13}
\loigiai{
\textbf{Đáp số:} $13$\\
$15 - 7 + 5 = 8 + 5 = 13$ viên bi.
}
\end{ex}

\begin{ex} %[2D2V13]
Một cửa hàng có 18 kg đường. Buổi sáng bán 6 kg, buổi chiều bán 5 kg. Hỏi cửa hàng còn lại bao nhiêu ki-lô-gam đường?
\shortans{7}
\loigiai{
\textbf{Đáp số:} $7$\\
$18 - 6 - 5 = 7\text{ kg}$.
}
\end{ex}

\begin{ex} %[2D2V13]
Trên sân có 9 con chim bồ câu, một đàn khác bay đến có 7 con. Lát sau có 4 con bay đi. Hỏi trên sân lúc này còn bao nhiêu con chim bồ câu?
\shortans{12}
\loigiai{
\textbf{Đáp số:} $12$\\
$9 + 7 - 4 = 12$ con.
}
\end{ex}
